\chapter{Nausea}

\section*{Definition}
Unease stomach vomit often precedes vomiting alone emesis.

\section*{Pathophysiology}
Brainstem mediated reflex triggered GI irritation vestibular stimulation increased intracranial pressure metabolic disturbances psychogenic factors.

\section*{Etiology}
\begin{itemize}
  \item Viral gastroenteritis intoxication
  \item Pregnancy morning sickness
  \item Medication side effects opioids chemotherapy
  \item Gastric outlet obstruction migration
  \item Migraine headache associated
\end{itemize}

\section*{Indications for Evaluation}
Seek care if:
\begin{itemize}
  \item Severe or worsening symptoms
  \item Symptoms lasting more than Prolonged unable retain fluids dehydration blood vomit high fever suspected obstruction
  \item Accompanied by other concerning signs
\end{itemize}

\section*{Related Symptoms}
\begin{itemize}
  \item Abdominal cramps loss appetite dehydration vomiting headache
\end{itemize}
\textbf{Note:} Always consider serious underlying conditions when evaluating persistent nausea.

\section*{Self-Care / Home Management}
\begin{itemize}
  \item Maintain hydration with clear fluids (water, oral rehydration solutions, clear broths) in small, frequent amounts.
  \item Eat bland, easy-to-digest foods (bananas, rice, applesauce, toast) and avoid fatty, spicy, or dairy-heavy meals during recovery.
  \item Rest the digestive system by eating smaller meals more frequently rather than large meals.
  \item Over-the-counter remedies may help: antacids for heartburn, loperamide for diarrhea, fiber supplements for constipation.
  \item Seek medical attention if symptoms persist beyond 48 hours, or if accompanied by severe pain, fever, or bloody stools.
\end{itemize}

\section*{Diagnostic Approach}
\begin{itemize}
  \item History covers onset, duration, stool frequency/character, associated symptoms (pain, fever, nausea), dietary triggers, and travel history.
  \item Physical examination includes abdominal auscultation, palpation for tenderness/masses, and assessment of hydration status.
  \item Stool studies (culture, ova/parasites, occult blood, fecal calprotectin) help identify infectious or inflammatory causes.
  \item Endoscopy or colonoscopy is indicated for chronic symptoms, unexplained weight loss, or concerning features.
\end{itemize}

\section*{Treatment / Management Overview}
\begin{itemize}
  \item Supportive care with hydration and dietary modification is first-line for most acute gastrointestinal symptoms.
  \item Probiotics may help restore gut flora balance after infectious diarrhea or antibiotic use.
  \item Specific pharmacotherapy depends on diagnosis: antiemetics for vomiting, antispasmodics for cramping, antibiotics for bacterial infection.
  \item Chronic conditions (IBS, IBD, GERD) require specialist management with tailored medical and lifestyle interventions.
\end{itemize}

\clearpage
